% GENERATED FILE. Edit canonical lessons, then run npm run generate:books.
% canonical-source-hash: fnv1a64:14e2cdb9f2597f07
% canonical-lessons: JA-C102-sentakumono, JA-C102-senzai, JA-C102-yunomi, JA-C102-kama, JA-C102-manaita

\chapter{Laundry, Detergent, Teacup, Rice pot, Chopping board}
\label{ch:ja-laundry-detergent-teacup-rice-pot-chopping-board}
\hlchaptermodality{102}

\begin{chapteropening}
\textbf{By the end of this chapter:} \emph{I can say the laundry, detergent, a teacup, a rice pot and a chopping board.}

Complete the last lesson of chapter 102: i can say the laundry, detergent, a teacup, a rice pot and a chopping board.
\end{chapteropening}

\section[\texorpdfstring{\ja{せんたくもの} (sentakumono)}{せんたくもの (sentakumono)}]{\ja{せんたくもの} (sentakumono) — the laundry}

\label{lesson:JA-C102-sentakumono}



\begin{quote}
Before the new one: say the Japanese for a heated table, then the Japanese for heating.
\end{quote}

\subsection*{You'll want to know: \ja{せんたくもの}}
\textbf{\ja{せんたくもの}} — \emph{sentakumono} — \textquotedblleft{}the laundry\textquotedblright{}.

\begin{grammarlens}[title={What you've built}]
One more word for this chapter.
\end{grammarlens}

\subsection*{Your turn}
\begin{itemize}
\raggedright
  \item \emph{Say it:} \emph{sentakumono}
  \item \emph{Say it:} \emph{sentakumono}, once more
  \item \emph{Say it:} say \emph{sentakumono}
  \item \emph{Recall:} say the Japanese for a heated table, then the Japanese for heating, then say \emph{sentakumono} again
\end{itemize}

\begin{tcolorbox}[breakable,colback=teal!4,colframe=teal!35!black,title={Before you move on}]
What does \ja{せんたくもの} mean? (\textquotedblleft{}The laundry\textquotedblright{}.) Say it once more.
\end{tcolorbox}
\section[\texorpdfstring{\ja{せんざい} (senzai)}{せんざい (senzai)}]{\ja{せんざい} (senzai) — detergent}

\label{lesson:JA-C102-senzai}



\begin{quote}
Before the new one: say the Japanese for heating, then the Japanese for the laundry.
\end{quote}

\subsection*{You'll want to know: \ja{せんざい}}
\textbf{\ja{せんざい}} — \emph{senzai} — \textquotedblleft{}detergent\textquotedblright{}.

\begin{grammarlens}[title={What you've built}]
One more word for this chapter.
\end{grammarlens}

\subsection*{Your turn}
\begin{itemize}
\raggedright
  \item \emph{Say it:} \emph{senzai}
  \item \emph{Say it:} \emph{senzai}, once more
  \item \emph{Say it:} say \emph{senzai}
  \item \emph{Recall:} say the Japanese for heating, then the Japanese for the laundry, then say \emph{senzai} again
\end{itemize}

\begin{tcolorbox}[breakable,colback=teal!4,colframe=teal!35!black,title={Before you move on}]
What does \ja{せんざい} mean? (\textquotedblleft{}Detergent\textquotedblright{}.) Say it once more.
\end{tcolorbox}
\section[\texorpdfstring{\ja{ゆのみ} (yunomi)}{ゆのみ (yunomi)}]{\ja{ゆのみ} (yunomi) — a teacup}

\label{lesson:JA-C102-yunomi}



\begin{quote}
Before the new one: say the Japanese for the laundry, then the Japanese for detergent.
\end{quote}

\subsection*{You'll want to know: \ja{ゆのみ}}
\textbf{\ja{ゆのみ}} — \emph{yunomi} — \textquotedblleft{}a teacup\textquotedblright{}.

\begin{grammarlens}[title={What you've built}]
One more word for this chapter.
\end{grammarlens}

\subsection*{Your turn}
\begin{itemize}
\raggedright
  \item \emph{Say it:} \emph{yunomi}
  \item \emph{Say it:} \emph{yunomi}, once more
  \item \emph{Say it:} say \emph{yunomi}
  \item \emph{Recall:} say the Japanese for the laundry, then the Japanese for detergent, then say \emph{yunomi} again
\end{itemize}

\begin{tcolorbox}[breakable,colback=teal!4,colframe=teal!35!black,title={Before you move on}]
What does \ja{ゆのみ} mean? (\textquotedblleft{}A teacup\textquotedblright{}.) Say it once more.
\end{tcolorbox}
\section[\texorpdfstring{\ja{かま} (kama)}{かま (kama)}]{\ja{かま} (kama) — a rice pot}

\label{lesson:JA-C102-kama}



\begin{quote}
Before the new one: say the Japanese for detergent, then the Japanese for a teacup.
\end{quote}

\subsection*{You'll want to know: \ja{かま}}
\textbf{\ja{かま}} — \emph{kama} — \textquotedblleft{}a rice pot\textquotedblright{}.

\begin{grammarlens}[title={What you've built}]
One more word for this chapter.
\end{grammarlens}

\subsection*{Your turn}
\begin{itemize}
\raggedright
  \item \emph{Say it:} \emph{kama}
  \item \emph{Say it:} \emph{kama}, once more
  \item \emph{Say it:} say \emph{kama}
  \item \emph{Recall:} say the Japanese for detergent, then the Japanese for a teacup, then say \emph{kama} again
\end{itemize}

\begin{tcolorbox}[breakable,colback=teal!4,colframe=teal!35!black,title={Before you move on}]
What does \ja{かま} mean? (\textquotedblleft{}A rice pot\textquotedblright{}.) Say it once more.
\end{tcolorbox}
\section[\texorpdfstring{\ja{まないた} (manaita)}{まないた (manaita)}]{\ja{まないた} (manaita) — a chopping board}

\label{lesson:JA-C102-manaita}



\begin{quote}
Before the new one: say the Japanese for a teacup, then the Japanese for a rice pot.
\end{quote}

\subsection*{You'll want to know: \ja{まないた}}
\textbf{\ja{まないた}} — \emph{manaita} — \textquotedblleft{}a chopping board\textquotedblright{}.

\begin{grammarlens}[title={What you've built}]
One more word for this chapter.
\end{grammarlens}

\subsection*{Your turn}
\begin{itemize}
\raggedright
  \item \emph{Say it:} \emph{manaita}
  \item \emph{Say it:} \emph{manaita}, once more
  \item \emph{Say it:} say \emph{manaita}
  \item \emph{Recall:} say the Japanese for a teacup, then the Japanese for a rice pot, then say \emph{manaita} again
\end{itemize}

\begin{tcolorbox}[breakable,colback=teal!4,colframe=teal!35!black,title={Before you move on}]
What does \ja{まないた} mean? (\textquotedblleft{}A chopping board\textquotedblright{}.) Say it once more.
\end{tcolorbox}
